\section{Imperative Syntax}\label{sec:syntax-design}

Motivo uses C-style imperative syntax---blocks, semicolons, and familiar control flow---so composers can focus on
musical logic rather than a novel notation paradigm.
\texttt{loop} and \texttt{for} express repetition; \texttt{if}/\texttt{else} express conditional variation; typed
variables let loop indices and parameters drive pitch or rhythm algorithmically.
Because the compiler unrolls control flow at build time, the composer gets procedural generation without a
real-time interpreter.

\begin{lstlisting}[language=Motivo, caption={A compact Motivo drum pattern.}, label={lst:motivo-groove}]
tempo 128;
signature 4/4;

track percussion using drums {
    pattern bar() {
        for (int i = 0; i < 16; i = i + 1) {
            play hihat from i;
            if (i % 8 == 0) { play kick from i; }
            if (i % 8 == 4) { play snare from i; }
        }
    }

    loop (500) { play bar(); }
}
\end{lstlisting}
